\textbf{Proof.}
An interval-valued procedure must be a measurable map from \(\mathcal D_{n,m}\) to subsets of \(\mathbb R\).  The adaptive handle only instructs the reader to use dyadic differences, estimate \(a\), combine diagnostics, and select a pair \((h,J)\).  It provides no formula for any diagnostic threshold, no map from their values to \((h,J)\), and no tie-breaking convention.  In particular, the deterministic rule that always chooses the smallest candidate bandwidth and the deterministic rule that always chooses the largest candidate bandwidth both satisfy the literal instruction to select a pair from the displayed family, yet they are different maps and return different intervals on samples for which those candidate intervals differ.  Thus the definition does not name one random interval.  Since both coverage and expected length are properties of a specified random interval, neither quantity is defined for this handle, and the requested positive implication cannot follow from the frozen core. \(\square\)